\section{Math-Shepherd: sin anotación humana de pasos}

\subsection{Definir la calidad del paso con "el potencial de continuar hasta la respuesta correcta"}

Math-Shepherd define la calidad de un prefijo de razonamiento como: la probabilidad de que, al continuar el muestreo a partir de este prefijo, se llegue finalmente a la respuesta correcta. Para el prefijo $s_{1:t}$ se muestrean $K$ rollouts subsecuentes:

$$
\hat q_t=\frac{1}{K}\sum_{k=1}^{K}
\mathbf 1\!\left[\mathrm{Answer}(s_{1:t},\tilde s^{(k)}_{t+1:T})=a^*\right].
$$

\begin{itemize}
    \item $\tilde s^{(k)}_{t+1:T}$: la $k$-ésima trayectoria subsecuente generada al continuar desde el prefijo actual;
    \item $a^*$: la respuesta final ground-truth;
    \item $\hat q_t$: la probabilidad empírica de que este prefijo de paso pueda conducir a la solución correcta.
\end{itemize}

\videofigure{figures/math-shepherd-annotation.jpg}{Math-Shepherd usa la tasa de éxito de los rollouts subsecuentes para etiquetar automáticamente la calidad del paso.}{00:38:30--00:40:42}

\subsection{Hard y Soft Annotation}

\begin{itemize}
    \item \textbf{Hard label}: en cuanto exista un rollout exitoso, se etiqueta el prefijo como positivo;
    \item \textbf{Soft label}: se usa la proporción de rollouts exitosos $\hat q_t$;
    \item El hard label enfatiza "si todavía hay esperanza", el soft label refleja "qué tanta esperanza hay".
\end{itemize}

\videofigure{figures/step-label-example.jpg}{Ejemplo de las etiquetas de outcome y de potencial de paso para un mismo proceso de solución.}{00:40:42--00:43:36}

\begin{knowledgebox}{Esto es un tipo de Monte Carlo value estimation}
Math-Shepherd en realidad está estimando el valor del estado: dado el prefijo de razonamiento actual, con qué probabilidad continuar ejecutando la política actual llevará a la respuesta correcta. No requiere que un humano juzgue la lógica de cada paso, solo requiere que la respuesta final pueda verificarse automáticamente.
\end{knowledgebox}

\subsection{Un verifier, tres formas de uso}

\videofigure{figures/math-shepherd-uses.jpg}{Un PRM entrenado automáticamente puede usarse para rerank, y también como PPO reward.}{00:43:36--00:44:36}

\begin{enumerate}
    \item \textbf{Verification}: ordenar múltiples respuestas completas;
    \item \textbf{Search}: puntuar y podar prefijos intermedios;
    \item \textbf{Reinforcement learning}: usar el reward por paso en PPO para mejorar directamente el generator.
\end{enumerate}

\videofigure{figures/math-shepherd-verification.jpg}{Math-Shepherd mejora la selección de candidatos en GSM8K y MATH500.}{00:44:36--00:47:02}

\videofigure{figures/math-shepherd-rl.jpg}{PPO con el PRM automático como reward mejora aún más al generator.}{00:47:02--00:48:06}

\subsection{Fuentes de sesgo en la anotación automática}

\begin{itemize}
    \item Cuando la rollout policy es demasiado débil, incluso un prefijo correcto puede no ser completado por nadie, y se etiqueta erróneamente como de bajo valor;
    \item Con un número limitado de rollouts, la varianza de $\hat q_t$ es muy grande;
    \item Que la respuesta final sea verificable no implica que el razonamiento intermedio sea realmente válido;
    \item Cuando cambian el generator de entrenamiento y el generator de los rollouts, el valor de los pasos sufre un desplazamiento de distribución.
\end{itemize}

\begin{warningbox}{"Sin anotación humana" no equivale a sin supervisión}
Math-Shepherd sigue dependiendo de una respuesta final verificable automáticamente, y consume una gran cantidad de compute de rollouts. Convierte el costo de mano de obra en costo de muestreo, y es aplicable a dominios verificables como matemáticas y código; para tareas abiertas, el reward final en sí mismo sigue siendo un problema difícil.
\end{warningbox}

\subsection{Resumen de la sección}

Math-Shepherd usa el resultado final verificable del entorno para inferir el valor de los pasos hacia atrás, logrando process supervision automática, y usa el verifier tanto para la inferencia como para el RL. Muestra el prototipo de un ciclo cerrado de automejora: el modelo genera trayectorias, la retroalimentación del resultado final etiqueta los pasos, y el PRM guía al modelo a generar mejores trayectorias.
